% !TEX root = chapter-standalone.tex

\section[Actions]{Actions}
\label{sec:cat-act}

In previous sections, we have seen how algebraic structures such as \SY{semigroups}, \SY{monoids}, and \SY{groups} form categories, where the morphisms are the appropriate structure-preserving maps (homomorphisms).
We now do the same for actions.

\subsection{The category of monoid actions}

We fix a \SY{monoid}~$\monoidA$ and organize all \SY{actions} of~$\monoidA$ into a category.

For this to work, we need that homomorphisms of \SY{monoid actions} compose. Let us check. 

Given homomorphisms~$\mapa \colon \tup{\setA, \act_\setA} \mto \tup{\setB, \act_\setB}$ and~$\mapb \colon \tup{\setB, \act_\setB} \mto \tup{\setC, \act_\setC}$, their composite~$\mapa \mthen \mapb \colon \setA \sto \setC$ is again a homomorphism of actions, since for every~$\monela \setin \monoidAset$,
\begin{equation}
    \act_\setA(\monela) \mthen (\mapa \mthen \mapb) = (\act_\setA(\monela) \mthen \mapa) \mthen \mapb = (\mapa \mthen \act_\setB(\monela)) \mthen \mapb = \mapa \mthen (\act_\setB(\monela) \mthen \mapb) = \mapa \mthen (\mapb \mthen \act_\setC(\monela)) = (\mapa \mthen \mapb) \mthen \act_\setC(\monela),
\end{equation}
using associativity of composition in~\Set and the equivariance conditions for~$\mapa$ and~$\mapb$.

Similarly, for any action~$\tup{\setA, \act_\setA}$, the identity function~$\catid_\setA \colon \setA \sto \setA$ is trivially a homomorphism of actions.

\begin{ctdefinition}[Category of $\monoidA$-actions]
    \SYNDEF{category of monoid actions}
    \label{def:cat-monoid-actions}
    Let~$\monoidA$ be a \SY{monoid}.
    The category of $\monoidA$-actions, denoted~$\monoidA\text{-}\Set$, is:
    \begin{enumerate}
        \item \emph{Objects}: all \SY{actions} of~$\monoidA$ on sets, i.e., pairs~$\tup{\setA, \act_\setA}$ where~$\setA$ is a set and~$\act_\setA \colon \monoidA \fto \Endof\setA$ is a \SY{monoid homomorphism}.
        \item \emph{Morphisms}: for \SY{actions}~$\tup{\setA, \act_\setA}$ and~$\tup{\setB, \act_\setB}$, the hom-set consists of all homomorphisms of $\monoidA$-actions (\cref{def:monoid-action-morphism}) from~$\tup{\setA, \act_\setA}$ to~$\tup{\setB, \act_\setB}$.
        \item \emph{Composition}: composition of the underlying functions.
        \item \emph{Identity morphisms}: for an action~$\tup{\setA, \act_\setA}$, its identity morphism is the identity function~$\catid_\setA$.
    \end{enumerate}
\end{ctdefinition}

\begin{remark}
    There is a \SY{forgetful functor}~$\monoidA\text{-}\Set \fto \Set$ which sends each action~$\tup{\setA, \act_\setA}$ to its underlying set~$\setA$, and each homomorphism of actions to the underlying function.
\end{remark}

\begin{remark}
    If we replace the \SY{monoid}~$\monoidA$ with a \SY{group}~$\grpA$, then the analogous construction yields a category of $\grpA$-actions (also called $\grpA$-sets), where the objects are \SY{group actions}~$\act \colon \grpA \fto \Autof\setA$ and the morphisms are equivariant maps satisfying the same condition \cref{eq:equivariance}.
\end{remark}
